\subsection{项目特色}
本项目将新疆光伏气象服务中的短时变化风险作为共同评价对象，在同一发布时间约束下检验可用数据、概率模型和事件预警。研究按发布时刻组织气象输入，分别报告正常运行、天气扰动和停限机条件，重视容易被总体平均误差掩盖的天气转换过程。依托已有多尺度时序研究和计算条件，控制模型数量，以小型基金能够承担的计算实验和离线原型形成可复核结果。

\subsection{创新点}
\subsubsection{天气变化条件下跨尺度预测增益与校准失配的联合辨析}
已有多尺度建模和光伏共形预测为本项目提供了方法起点。本项目进一步在相同可用信息和计算预算下，考察气象输入、尺度选择及历史误差更新之间的关系，检验某些天气中的精度增益能否与概率可靠性同时保持。预期贡献是识别稳定增益与失配的条件，给出可复核的变量、尺度和校准选择依据；将现有算法并置不作为创新证据。

\subsubsection{面向爬坡预警的时间依赖贡献及失效条件检验}
本项目保持各时距预测样本的边际一致，通过保留或打散历史秩依赖构造受控比较，检验天气转换前后的依赖迁移是否影响爬坡概率。与同时更换预测器和场景生成器的比较相比，该设计能够把事件概率变化与依赖结构对应起来，并用概率评分、事件级告警及反证实验判断方法何时有效。秩重排采用已有方法，拟形成的新认识是边际固定时，天气条件历史秩相对总体历史秩对事件可靠性的增量及失效边界。该认识限定为目标场景的经验适用条件；天气条件选择无效或造成反向影响也构成应报告的结果。区间校准与事件预测分别评价，不以区间覆盖改善证明事件概率随之改善。
